Classical statistical learning typically assumes a stationary data-generating law and a fixed set of observed variables. Many practical applications, however, depart from this baseline. Strategic agents manipulate or selectively disclose information in response to deployed decision rules, while in other domains, features that are jointly predictive are never observed together during training. This dissertation investigates these two distinct structural departures.

Within the strategic framework, we examine three problems: strategic classification under non-uniform preferences, strategic feature revelation in regression, and strategic window dressing in Initial Public Offering (IPO) regulation. Formulating these settings as Stackelberg screening games, we derive structural properties of optimal decision rules and develop learning algorithms with provable generalization and regret guarantees. Notably, we show that accounting for strategic incentives not only provides robustness against manipulation, but can also elicit valuable private information that strictly improves performance over non-strategic baselines. In IPO regulation, we show that social welfare can attain a strict interior optimum, revealing an inherent regulatory trade-off.

For partitioned observations, which represent a structural rather than strategic departure, predictive features and labels are partitioned across training sets and never jointly observed. We characterize the resulting learning limits, develop an imputation-based empirical risk minimization framework, and identify conditions under which the optimal regularizer follows an interpretable three-candidate rule.

Together, these results show that data-driven decision making depends fundamentally on agent incentives and observation structure alongside the data-generating law, extending statistical learning theory to richer operational environments.